\section{Relative tails and computable hopping calibration}
\label{sec:hopping}

The absolute orbital estimate of Proposition~\ref{prop:isolated-well}
does not determine an overlap of order $e^{-D}$.  This section proves a
relative tail estimate, identifies positive integrals that determine the
hoppings, and shows how to calibrate those integrals before tuning the
single-well energies.  All constants below are independent of the site
layout, $D$, the displacements, and the parameters $\rho_i\in[1,4]$.
The cap profile $\chi$ is fixed throughout.

We use the positions and operators from Section~\ref{sec:wells}.  In
particular,
\[
 X_i=Dp_i+s_i,\qquad \max_i|s_i|\le S,
 \qquad |X_i-X_j|\ge D-2S\quad(i\ne j).
\]
We take $D$ sufficiently large and, for the estimates below, may assume
$m+S\le D^{1/4}$.  Where a relative error is expanded, its displayed
bound is assumed smaller than a fixed positive constant.  The final
parameter choice in Section~\ref{sec:completion} satisfies these
conditions and makes every such error as small as required.

\subsection{Relative comparison of single-well tails}

For $r_i=|y-X_i|$ define the ray primitive
\begin{equation}
 F_i(y)=r_i\int_0^1 W_i\bigl(X_i+t(y-X_i)\bigr)\,\dd t,
 \qquad F_i(X_i)=0.
 \label{eq:ray-primitive}
\end{equation}
The integral follows the entire segment from $X_i$ to $y$; it is not
a local approximation to the field at either endpoint.  Write
\begin{equation}
 \varepsilon_{\mathrm{rel}}
   =\frac{m\log D+m^2}{D}.
 \label{eq:relative-tail-error}
\end{equation}

\begin{proposition}[Uniform relative tail]
\label{prop:relative-tail}
Suppose the actual rounded charges and the parameters $\rho_i$ obey
the conclusions of Proposition~\ref{prop:tuning}, in particular
$|\lambda_i+1/2|\le e^{-1.01D}$.  Uniformly for $r_i\le2D$,
\begin{equation}
 \psi_i(y)=\pi^{-1/2}e^{-r_i+F_i(y)}
       \bigl(1+O(\varepsilon_{\mathrm{rel}})\bigr).
 \label{eq:relative-tail}
\end{equation}
\end{proposition}

\begin{proof}
Fix a site, suppress its subscript, and write $r=|y-X|$ and
$\omega=(y-X)/r$ when $r>0$.  The capped-field estimates in
Proposition~\ref{prop:isolated-well} give
\begin{equation}
 \|W\|_\infty\le \frac{Cm}{D},\qquad
 \|\nabla W\|_\infty\le \frac{Cm}{D^2},\qquad
 \|D^2W\|_\infty\le \frac{Cm}{D^3}.
 \label{eq:hopping-field-derivatives}
\end{equation}
Set $A(y)=\int_0^1W(X+t(y-X))\,\dd t$, so $F=rA$.
For $0<r\le3D$, differentiation gives
\begin{align*}
 \nabla F
  &=\omega A+r\int_0^1t\nabla W(X+t(y-X))\,\dd t,\\
 \Delta F
  &=\frac{2A}{r}
    +2\omega\cdot\int_0^1t\nabla W(X+t(y-X))\,\dd t\\
  &\quad+r\int_0^1t^2\Delta W(X+t(y-X))\,\dd t.
\end{align*}
Consequently,
\begin{equation}
 |\nabla F|\le \frac{Cm}{D},\qquad
 |\Delta F|\le C\left(\frac{m}{Dr}+\frac{m}{D^2}\right).
 \label{eq:ray-derivative-bounds}
\end{equation}
Also $F(r,\omega)=\int_0^rW(X+t\omega)\,\dd t$, and hence
$\partial_rF=W(y)$ exactly.

Let $f=\pi^{-1/2}e^{-r+F}$.  A direct calculation on the annulus
$10\le r\le3D$ gives
\begin{align}
 \frac{(h_i-\lambda_i)f}{f}
  &=-(\lambda_i+1/2)+\frac1r-V_i
       -\frac12\bigl(\Delta F+|\nabla F|^2\bigr),
       \label{eq:relative-tail-residual}\\
 \left|\frac{(h_i-\lambda_i)f}{f}\right|
  &\le C\left(\frac{m}{Dr}+\frac{m^2}{D^2}\right).
       \label{eq:relative-tail-residual-bound}
\end{align}
In the first equality the term $\partial_rF$ cancels $W_i$.
For the second, note that $\rho_i\le4$ and $r\ge10$, so
\[
 \left|\frac1r-V_i(y)\right|
 \le \frac{|1-q|}{r}+\frac{q'}{r-4}
 \le \frac{Cm}{Dr}.
\]
The charge-rounding errors and $e^{-1.01D}$ are absorbed by the
right-hand side of \eqref{eq:relative-tail-residual-bound}.

Choose a sufficiently large fixed constant $A_0$ and put
\begin{equation}
 g(r)=A_0\left[
     \frac{m}{D}\bigl(1+\log(r/10)\bigr)
     +\frac{m^2r}{D^2}\right].
 \label{eq:tail-barrier-correction}
\end{equation}
Then
\[
 g'=A_0\left(\frac{m}{Dr}+\frac{m^2}{D^2}\right),\qquad
 \Delta g=A_0\left(\frac{m}{Dr^2}+\frac{2m^2}{D^2r}\right),
 \qquad 0\le\Delta g\le\frac{2g'}r.
\]
Replacing $f$ by $fe^{\pm g}$ adds to its operator quotient the term
\begin{equation}
 \pm(1-\partial_rF)g'
       \mp\frac12\Delta g-\frac12(g')^2.
 \label{eq:barrier-signs}
\end{equation}
Since $r\ge10$, $\partial_rF=O(m/D)$, and $g'$ is uniformly small
for sufficiently large $D$, the added term is at least $c g'$ for
the plus sign and at most $-c g'$ for the minus sign.  Increasing
$A_0$ dominates \eqref{eq:relative-tail-residual-bound}.  Thus
$fe^g$ is a supersolution and $fe^{-g}$ a subsolution.

At $r=10$, Proposition~\ref{prop:isolated-well} and Sobolev embedding
give $\psi_i=\pi^{-1/2}e^{-10}+O(m/D)$, while $F=O(m/D)$.
The reference value $e^{-10}$ is a fixed positive number.  A further
fixed increase of $A_0$ therefore ensures
$fe^{-g}\le\psi_i\le fe^g$ on the inner boundary.
On the outer boundary $r=3D$, Lemma~\ref{lem:well-tail} gives
\[
 |\psi_i|\le C e^{-0.99(3-0.01)D}=C e^{-2.9601D}.
\]
Moreover $0\le F\le Cm$ there.  With $\kappa=e^{-2.5D}$, the
functions
\[
 f_+=fe^g+\kappa,\qquad f_-=fe^{-g}-\kappa
\]
satisfy $f_-\le\psi_i\le f_+$ on both boundary spheres for large
$D$.  On the annulus the zero-order coefficient obeys
\[
 -V_i-W_i-\lambda_i
 \ge \frac12-\frac1{10}-O(m/D)\ge\frac14.
\]
Adding the positive constant $\kappa$ preserves the supersolution
inequality, and subtracting it preserves the subsolution inequality.
The maximum principle therefore yields the same inequalities throughout
the annulus.  For completeness, its weak form here follows by testing
a supersolution-minus-solution with its nonnegative negative part:
the resulting identity is the negative of a sum of a gradient norm
and a strictly positive weighted $L^2$ norm, so the negative part
vanishes.  The same argument applies to a solution-minus-subsolution.

On $10\le r\le2D$, one has
$g(r)=O(\varepsilon_{\mathrm{rel}})$ and, since $F\ge0$,
\[
 \frac{\kappa}{f}\le \sqrt\pi\,e^{-D/2}.
\]
The comparison inequalities imply \eqref{eq:relative-tail} on this
region.  On $r\le10$, the reference exponential is bounded below by
a fixed positive constant, $F=O(m/D)$, and the uniform absolute
orbital estimate gives the same relative conclusion directly.
\end{proof}

\subsection{Surrogate fields and positive hopping integrals}

We first identify the attraction integral to be calibrated. Write
$V_k^{\mathrm{hole}}=(1-\chi_k)V_k$ for each site $k$, and for an
ordered pair of distinct sites $(i,j)$ define
\begin{equation}
 B_{ij}=\langle\psi_j,V_j^{\mathrm{hole}}\psi_i\rangle.
 \label{eq:oriented-hole-integral}
\end{equation}
The identity
\[
 h=h_i-\sum_{k\ne i}V_k^{\mathrm{hole}}
\]
shows that the hole at $j$ contributes $-B_{ij}$ to the matrix
element $\langle\psi_j,(h+1/2)\psi_i\rangle$ in the nonorthogonal
orbitals. Section~\ref{sec:continuum} will bound its remaining terms
and the change under orthonormalization.

The relative tail formula still contains the secondary parameters, which
have not yet been computed when the primary centers must be placed. We
replace its capped field by an explicit one depending only on the primary
centers. The replacement error is relative and uniform in the entire
tuning cube; this is what makes the construction order possible.

Define
\begin{equation}
 W_i^*(y)=\sum_{k\ne i}\frac{\chi(|y-X_k|/D)}{|y-X_k|},\qquad
 F_i^*(y)=r_i\int_0^1W_i^*(X_i+t(y-X_i))\,\dd t.
 \label{eq:surrogate-field}
\end{equation}
Each summand is extended by zero through its cutoff hole.  In
particular the apparent quotient at $y=X_k$ is identically zero
in a neighborhood of that point.

\begin{lemma}[Uniform capped kernels and surrogate replacement]
\label{lem:surrogate}
The zero-extended kernels
\[
 G_D(z)=\frac{\chi(|z|/D)}{|z|},\qquad
 J_{D,\rho}(z)=\frac{\chi(|z|/D)}{|z-\rho e_z|}
\]
are globally $C^2$, uniformly for $1\le\rho\le4$.  For every
multi-index $\alpha$ with $|\alpha|\le2$,
\begin{equation}
 \|\partial^\alpha G_D\|_\infty+
 \|\partial^\alpha J_{D,\rho}\|_\infty
 \le C D^{-1-|\alpha|}.
 \label{eq:capped-kernel-bounds}
\end{equation}
Uniformly over all allowed secondary positions and rounded charges,
\begin{equation}
 \|W_i-W_i^*\|_\infty\le\frac{Cm^2}{D^2},\qquad
 |F_i(y)-F_i^*(y)|\le\frac{Cm^2}{D}\quad(r_i\le3D).
 \label{eq:surrogate-error}
\end{equation}
Consequently, for tuned wells the relative formula
\eqref{eq:relative-tail} holds with $F_i^*$ in place of $F_i$ and
the same error bound $O(\varepsilon_{\mathrm{rel}})$.
\end{lemma}

\begin{proof}
For large $D$, both poles are strictly inside $|z|<D/16$, where
the kernels vanish.  On the nonzero part of either kernel,
$|z|\ge D/16$ and $|z-\rho e_z|\ge |z|/2$.  Rescale by $z=D\zeta$.
Both kernels are $D^{-1}$ times a function whose derivatives through
order two are uniformly bounded: the shifted pole in the second
rescaled kernel has position $(\rho/D)e_z$, a fixed positive
distance from the support of $\chi$.  The $C^2$ matching of $\chi$
at its transition endpoints makes these global bounds, including
at the hole boundaries.  This proves \eqref{eq:capped-kernel-bounds}.

Since $|q-1|+q'\le Cm/D$ after rounding,
\[
 W_i-W_i^*=
  \sum_{k\ne i}\bigl[(q-1)G_D(y-X_k)
                  +q'J_{D,\rho_k}(y-X_k)\bigr]
\]
has sup norm $O(m^2/D^2)$.  Integration over a ray of length at
most $3D$ proves the second bound in \eqref{eq:surrogate-error}.
Exponentiating this bound in Proposition~\ref{prop:relative-tail}
proves the final assertion.  None of the kernel estimates require
tuning; they hold uniformly on the whole parameter cube.
\end{proof}

For an oriented contact $(i,j)$, replacing the two orbital factors
in $B_{ij}$ by their relative
approximations and the local nuclear potential by the unit primary
kernel leads to the explicit positive integral
\begin{align}
 T_{ij}(X)
  &=\frac1\pi\int_{|v|\le D/8}
       \frac{1-\chi(|v|/D)}{|v|}\,e^{L_X(v)}\,\dd^3v,
       \label{eq:positive-hopping-integral}\\
 L_X(v)
  &=-|X_j+v-X_i|-|v|
     +F_i^*(X_j+v)+F_j^*(X_j+v).
       \label{eq:hopping-log-integrand}
\end{align}
The fields in this definition use all centers in $X$.  The
integrable singularity $|v|^{-1}$ is part of the nonnegative
integration weight, not part of the exponential. The next lemma
justifies these replacements.

\begin{lemma}[Comparison with the uncapped attraction]
\label{lem:positive-hopping}
For tuned wells and every contact,
\begin{equation}
 B_{ij}=T_{ij}(X)
       \bigl(1+O(\varepsilon_{\mathrm{rel}})\bigr).
 \label{eq:hole-surrogate-comparison}
\end{equation}
The constants are uniform in $\rho_j\in[1,4]$.
\end{lemma}

\begin{proof}
On the support of $V_j^{\mathrm{hole}}$, write $y=X_j+v$.
Then $|v|\le D/8$ and
$|y-X_i|\le D+2S+D/8<2D$.  Proposition~\ref{prop:relative-tail}
and Lemma~\ref{lem:surrogate} apply to both wave functions.  Their
product is $\pi^{-1}e^{L_X(v)}$ times a uniform
$1+O(\varepsilon_{\mathrm{rel}})$ factor.  The primary kernel has
coefficient $q=1+O(m/D)$, so it remains to compare the secondary
kernel with the primary one in the resulting positive integral.

For $|v|>10$,
\[
 \frac{|v|}{|v-\rho_j e_z|}
   \le\frac{|v|}{|v|-4}\le\frac53.
\]
For $|v|\le10$, the hole weight is one when $D$ is large.  The
ray derivative estimates, also valid for $F^*$, show that $L_X$
is Lipschitz on this ball with a fixed constant.  This statement
requires no differentiability of $|v|$ at zero; it follows either
by integrating the bounded derivatives away from zero or directly
from the ray formula.  Thus $e^{L_X(v)}$ is within fixed factors
of $e^{L_X(0)}$ throughout the ball.  Moreover,
\[
 \int_{|v|\le10}\frac{\dd^3v}{|v-\rho_j e_z|}
  \le\int_{|w|\le14}\frac{\dd^3w}{|w|}<\infty,
 \qquad
 \int_{|v|\le1}\frac{\dd^3v}{|v|}>0.
\]
The secondary-kernel integral on the inner ball is therefore at
most a fixed multiple of the primary-kernel integral.  Combining
the two regions and multiplying by $q'=O(m/D)$ gives a relative
$O(m/D)$ correction.  Positivity allows the uniform relative
wave-function errors to be taken outside both integrals.  These
errors are included in \eqref{eq:hole-surrogate-comparison}.
\end{proof}

\begin{lemma}[Size of the reference integrals]
\label{lem:reference-hopping}
Let $X^0=(Dp_i)_i$.  For every contact,
\begin{equation}
 c e^{-D}\le T_{ij}(X^0)\le C D^2 e^{-D+Cm}.
 \label{eq:reference-hopping-bounds}
\end{equation}
\end{lemma}

\begin{proof}
Put $u=u_{ij}=p_j-p_i$, so $|u|=1$.  The triangle inequality
gives $|Du+v|+|v|\ge D$.  Each ray correction in the integral is
nonnegative and at most $Cm$, by \eqref{eq:capped-kernel-bounds}
and its ray length.  The integral of $|v|^{-1}$ over the ball of
radius $D/8$ is $O(D^2)$, proving the upper bound.  On $|v|\le1$
the hole weight is one, $|Du+v|+|v|\le D+2$, and the corrections
are nonnegative.  Integration over this fixed ball gives the lower
bound; for example $2e^{-D-2}$ is a valid lower bound for large $D$.
\end{proof}

\Needspace{10\baselineskip}
\subsection{Sensitivity to displacement}

\begin{lemma}[Uniform displacement law]
\label{lem:hopping-displacement}
Let $X_i=Dp_i+s_i$ with $|s_i|\le S$.  There is a constant
$c_1>0$ such that, uniformly on all contacts,
\begin{align}
 \frac{T_{ij}(X)}{T_{ij}(X^0)}
  &=e^{-u_{ij}\cdot(s_j-s_i)}
     \bigl(1+O(\varepsilon_{\mathrm{disp}})\bigr),
     \label{eq:hopping-displacement-law}\\
 \varepsilon_{\mathrm{disp}}
  &=\frac{mS}{D}+\frac{S^2}{D}
      +S D^{-1/4}+e^{-c_1\sqrt D}.
     \label{eq:hopping-displacement-error}
\end{align}
In particular the coefficient of the projected displacement in the
leading exponent is one.
\end{lemma}

\begin{proof}
Compare the two integrals at the same $v$.  In the ray defining
$F_i^*(X_j+v)$, the reference path point
$Dp_i+t(Dp_j+v-Dp_i)$ moves by $(1-t)s_i+ts_j$, of norm at
most $S$.  Each field center moves by at most $S$ as well.
The globally Lipschitz bound in \eqref{eq:capped-kernel-bounds}
therefore changes the field along the path by at most $CmS/D^2$.
The ray has length $O(D)$ and its length changes by at most $2S$.
Since the field itself is $O(m/D)$,
\begin{equation}
 |F_i^*(X_j+v;X)-F_i^*(Dp_j+v;X^0)|
    \le\frac{CmS}{D}.
 \label{eq:ray-displacement}
\end{equation}
The same bound holds for the $F_j^*$ term; its ray length $|v|$
does not change.  The global nature of the capped-kernel bound is
essential here: it also applies when a moving ray crosses a hole
or a transition interface.

Write $u=u_{ij}$, $\delta=s_j-s_i$, and $v=zu+w$ with
$w\perp u$.  Since $|v|\le D/8$,
$|Du+v|\ge7D/8$.  The Hessian of the norm on the segment from
$Du+v$ to $Du+v+\delta$ has norm $O(1/D)$.  Its gradient at
$Du+v$ differs from $u$ by $O(|w|/D)$.  Taylor's formula gives
\begin{equation}
 |Du+v+\delta|-|Du+v|
   =u\cdot\delta+O\left(\frac{S|w|}{D}+\frac{S^2}{D}\right).
 \label{eq:distance-displacement}
\end{equation}
This estimate includes the whole backward segment $v=-tu$,
$0\le t\le D/8$: there $Du+v=(D-t)u$ and the gradient is
exactly $u$.  Thus the segment's non-negligible longitudinal
extent creates neither a different coefficient nor an $O(S)$
transverse error.

On the region $|w|\le D^{3/4}$, Equations
\eqref{eq:ray-displacement} and \eqref{eq:distance-displacement}
show that
\begin{equation}
 L_X(v)-L_{X^0}(v)+u\cdot\delta
   =O\left(\frac{mS}{D}+\frac{S^2}{D}+S D^{-1/4}\right).
 \label{eq:main-region-log-error}
\end{equation}
The nonnegative integration weight is identical in the two
integrals, so exponentiation controls their contributions on
this region by the same relative factor.

On the remaining region $|w|>D^{3/4}$, one has
\begin{equation}
 \begin{split}
 |Du+v|+|v|-D
   &\ge z+|v|
    =\frac{|w|^2}{|v|-z}\\
   &\ge\frac{4|w|^2}{D}
    \ge4\sqrt D.
 \end{split}
 \label{eq:transverse-excess}
\end{equation}
The identity is used only for $w\ne0$, and the denominator is
at most $2|v|\le D/4$.  The field corrections are bounded by
$Cm$, without any angular constancy assumption.  Displacement
changes a distance by at most $2S$, and $|u\cdot\delta|\le2S$.
It follows that the omitted reference and displaced contributions,
after division by $e^{-D}$ and $e^{-D-u\cdot\delta}$ respectively,
are at most
\begin{equation}
 C D^2\exp\bigl[-c\sqrt D+C(m+S)\bigr].
 \label{eq:transverse-integral-error}
\end{equation}
By \eqref{eq:reference-hopping-bounds}, these are also bounds on
their contributions relative to the reference integral, with
the appropriate displacement factor.  The assumption
$m+S\le D^{1/4}$ and sufficiently large $D$ absorb both the
possibly large anisotropic factor $e^{Cm}$ and the polynomial
volume factor into $e^{-c_1\sqrt D}$.  Combining this with
\eqref{eq:main-region-log-error}, and using smallness of its
right-hand side, proves the lemma.
\end{proof}

\subsection{Deterministic relative quadrature}

We have reduced hopping calibration to positive integrals whose size is
exponentially small. Direct absolute quadrature at that scale would
obscure the complexity bound. Instead we control the oscillation of the
logarithm of the exponential factor, while integrating its nonnegative
radial weight exactly. Positivity then converts local errors into a
relative error for the whole integral.

\begin{lemma}[Computing the positive reference integrals]
\label{lem:hopping-quadrature}
For $0<\epsilon<1/10$, one can compute a positive rational
$\widehat I_{ij}$ satisfying
\begin{equation}
 \left|\log\frac{\widehat I_{ij}}{e^D T_{ij}(X^0)}\right|
   \le\epsilon
 \label{eq:reference-quadrature-accuracy}
\end{equation}
by a deterministic algorithm polynomial in $D$, $m$, the bit
lengths of the layout data, and $\epsilon^{-1}$.  In particular
the algorithm is polynomial in the source size when $D$ is a
fixed polynomial and $\epsilon$ an inverse polynomial.
\end{lemma}

\begin{proof}
Let $L=L_{X^0}$ in \eqref{eq:hopping-log-integrand}.  On the
integration ball, the two distance terms have Lipschitz constants
one, and the ray formulas and \eqref{eq:capped-kernel-bounds} give
a Lipschitz constant $Cm/D$ for each correction.  Hence
\begin{equation}
 |L(v)-L(v')|\le C|v-v'|,
 \qquad |L(v)|\le C(D+m).
 \label{eq:hopping-log-regularity}
\end{equation}
Only the exponential factor has been included in this logarithm.
No lower bound on $1-\chi$ or differentiability of $\log(1-\chi)$
is needed.

We describe a mesh with exact positive weights.  Parameterize
the unit sphere, up to a set of area zero, by
\[
 \omega(\mu,\varphi)
   =\bigl(\sqrt{1-\mu^2}\cos\varphi,
           \sqrt{1-\mu^2}\sin\varphi,\mu\bigr).
\]
Here $-1\le\mu\le1$, $0\le\varphi\le2\pi$, and
$\dd\omega=\dd\mu\,\dd\varphi$.  Use $M$ equal intervals
in $\mu$, $N$ equal intervals in $\varphi$, and a rational radial
partition with mesh at most $h$.  Insert all the finitely many
breakpoints of $\chi(r/D)$ into that partition.  On a radial
interval $I_a$, the weight
\[
 w_a=\int_{I_a}(1-\chi(r/D))r\,\dd r
\]
is a nonnegative exact rational, computable by polynomial
integration.  Discard intervals with zero weight.  The normalized
angular weight of every angular rectangle is exactly
\begin{equation}
 \frac1\pi\,\frac2M\,\frac{2\pi}{N}=\frac4{MN}.
 \label{eq:exact-angular-weight}
\end{equation}
Choose one sample $v_{abc}$ in each resulting spatial cell, for
example at the parameter midpoints.

The square-root parameterization satisfies
\[
 |\omega(\mu,\varphi)-\omega(\mu',\varphi')|
  \le C|\mu-\mu'|^{1/2}+C|\varphi-\varphi'|.
\]
Therefore $M=O((D/h)^2)$ and $N=O(D/h)$, with sufficiently
large fixed constants, make every spatial cell have diameter
$O(h)$, including cells adjacent to the poles.  There are
$O((D/h)^4)$ cells, a polynomial number.  Their exact weighted
sample sum is
\[
 Q=\sum_{a,b,c}\frac{4w_a}{MN}\,e^{L(v_{abc})}.
\]
Positivity and \eqref{eq:hopping-log-regularity} imply
\begin{equation}
 e^{-Ch}Q\le T_{ij}(X^0)\le e^{Ch}Q.
 \label{eq:positive-quadrature-bracket}
\end{equation}
Thus a fixed multiple of $\epsilon$ is sufficient for $h$.
This relative estimate is independent of the size $e^{-D}$ of
the integral.

It remains to specify evaluation and bit accuracy.  For a ray of
length $r\le2D$, the integrand
$t\mapsto W_i^*(X_i+t(y-X_i))$ has Lipschitz constant
at most $Cmr/D^2$.  A midpoint rule with $K$ equal panels
therefore approximates its ray primitive with error at most
\begin{equation}
 C\frac{mr^2}{D^2K}\le\frac{Cm}{K}.
 \label{eq:ray-quadrature-error}
\end{equation}
Taking $K=O(m/\epsilon)$ and evaluating each field sample to
absolute error $O(\epsilon/D)$ gives an $O(\epsilon)$ error
in each primitive.  Each field sample is a sum of $m-1$
globally Lipschitz capped kernels.  At small radius the kernel
is evaluated as zero, and elsewhere its denominator is bounded
below by a fixed multiple of $D$; equivalently one evaluates the
fixed continuous function $\chi(a)/a$, extended by zero on
$[0,1/16]$, at $a=|z|/D$.  Its global Lipschitz bound controls
errors even if the approximate radius lies on the other side of
a transition breakpoint.  No crossing location needs to be found.

The coordinates $Dp_i$, sample coordinates, square roots, and
trigonometric values can be computed to inverse-polynomial
absolute accuracy with polynomially many bits.  The global
kernel and ray bounds just proved control the resulting errors.
Thus all sample logarithms can be approximated to an arbitrarily
specified $O(\epsilon)$ absolute error with polynomial work.
Such logarithm errors change the positive sample sum by a factor
$e^{O(\epsilon)}$, not by an additive error multiplied by
an uncontrolled exponential.

For the remaining arithmetic, compute the scaled sum $e^D Q$,
using sample exponentials $e^{D+L(v_{abc})}$.  Its exact value
is bounded below by a positive constant up to the factor in
\eqref{eq:positive-quadrature-bracket}, because
\eqref{eq:reference-hopping-bounds} gives
$e^D T_{ij}(X^0)\ge2e^{-2}$.  The sum of its non-exponential
weights is $O(D^2)$, and all exponent arguments have magnitude
$O(D+m)$.  After the logarithms have been approximated as above,
each elementary exponential can be rounded to sufficiently small
absolute error $O(\epsilon/D^2)$, clipping a negative rounded
value to zero; this adds $O(\epsilon)$ to the scaled sum.  These
last approximations require only
$O(D+m+\log(D/\epsilon))$ precision bits, with polynomial
working precision.  This step does not require exponentially
accurate ray quadrature: the earlier ray errors have already
been bounded multiplicatively.  Standard range reduction and
Taylor series compute the square roots, trigonometric values,
exponentials, and logarithms to these precisions in polynomial
time.  The exact rational weights have polynomial bit lengths,
and the sum contains polynomially many terms.

Allocate a sufficiently small fixed fraction of $\epsilon$ to
each of the spatial, ray, coordinate, and arithmetic errors.
The resulting positive rational sum is $\widehat I_{ij}$ and
satisfies \eqref{eq:reference-quadrature-accuracy}.
\end{proof}

\subsection{Calibration, rounding, and orientation}

The oriented attraction integrals are close to a symmetric matrix
for an exact reason that is useful independently of their asymptotic
formula.

\begin{lemma}[Orientation identity]
\label{lem:hopping-orientation}
For distinct sites,
\begin{equation}
 B_{ij}-B_{ji}
   =(\lambda_i-\lambda_j)\langle\psi_j,\psi_i\rangle.
 \label{eq:hopping-orientation-identity}
\end{equation}
For tuned wells and large $D$ under the stated separation bounds,
\begin{equation}
 |\langle\psi_j,\psi_i\rangle|\le e^{-0.85D},\qquad
 |B_{ij}-B_{ji}|\le2e^{-1.86D}.
 \label{eq:hopping-orientation-error}
\end{equation}
\end{lemma}

\begin{proof}
The single-well operators have common domain $H^2$, and
\[
 h_i-h_j=V_j^{\mathrm{hole}}-V_i^{\mathrm{hole}}.
\]
Pair this identity between the real eigenfunctions $\psi_j$ and
$\psi_i$, using self-adjointness of $h_j$, to obtain
\eqref{eq:hopping-orientation-identity}.  All the pairings are
well-defined by the Coulomb domain bounds of
Proposition~\ref{prop:isolated-well}.  Equivalently, in the two
expansions of the full matrix element
$\langle\psi_j,(h+1/2)\psi_i\rangle$, contributions from every
hole other than $i$ and $j$ cancel identically, because they
integrate the same real product against the same multiplication
potential.

To record an explicit overlap bound, Lemma~\ref{lem:well-tail}
gives
\[
 |\psi_i(y)\psi_j(y)|
    \le C e^{0.0198D}e^{-0.99(r_i+r_j)}.
\]
Use $r_i+r_j\ge D-2S$ on a factor with exponent $0.90$, and
retain $e^{-0.09r_i}$ for integrability.  It follows that
\[
 |\langle\psi_j,\psi_i\rangle|
   \le C e^{-0.8802D+1.8S}
          \int_{\mathbb R^3}e^{-0.09|y-X_i|}\,\dd^3y
   \le e^{-0.85D}
\]
for sufficiently large $D$.  Finally
$|\lambda_i-\lambda_j|\le2e^{-1.01D}$ proves the second bound.
\end{proof}

For a tolerance $\tau>0$, define the total calibration error
\begin{equation}
 \begin{split}
 \mathcal E_{\mathrm{cal}}
  ={}&\frac{m\log D+m^2}{D}
       +\frac{mS}{D}+\frac{S^2}{D}\\
     &+S D^{-1/4}+\tau+e^{-c_1\sqrt D}.
 \end{split}
 \label{eq:total-calibration-error}
\end{equation}

\begin{proposition}[Deterministic hopping calibration]
\label{prop:calibration}
Let the contact layout be as in Lemma~\ref{lem:contact-layout},
choose one orientation per contact, and let the prescribed rational
hoppings satisfy $0<t_{\min}\le t_{ij}\le1$, where $t_{\min}$ is
an inverse polynomial in the source size.  Let $Z_*$ be an integer
with $e^D/2\le Z_*\le2e^D$ and round the nuclear charges as in
Section~\ref{sec:wells}.  For inverse-polynomial $\tau$, there is
a deterministic polynomial-time procedure that computes rational
centers $X_i=Dp_i+s_i$ obeying
\begin{equation}
 S:=\max_i|s_i|
   \le Cm\bigl(m+\log D+\log(1/t_{\min})+1\bigr)+\tau.
 \label{eq:calibrated-displacement-size}
\end{equation}
After applying Proposition~\ref{prop:tuning} at these final centers,
the resulting orbitals satisfy, on every chosen oriented contact,
\begin{equation}
 |Z_*B_{ij}-t_{ij}|
      \le C t_{ij}\mathcal E_{\mathrm{cal}}.
 \label{eq:calibrated-contact-error}
\end{equation}
Here $D$ is taken sufficiently large for the hypotheses above and
for $\mathcal E_{\mathrm{cal}}$ to be small.  The reverse oriented
attraction has an additional error at most $4e^{-0.86D}$ in the
same rescaled units.  No additional geometric equation for the
reverse orientation is needed.
\end{proposition}

\begin{proof}
Compute, to additive error at most $\tau/4$, the quantities
\begin{equation}
 d_{ij}=\log\left(\frac{Z_*T_{ij}(X^0)}{t_{ij}}\right),
 \qquad X^0=(Dp_i)_i.
 \label{eq:calibration-increments}
\end{equation}
Lemma~\ref{lem:hopping-quadrature} supplies the requisite relative
integral accuracy.  One can evaluate the logarithm as
\[
 \log Z_*-D+\log\bigl(e^D T_{ij}(X^0)\bigr)-\log t_{ij},
\]
so all elementary-function calculations have polynomial bit cost.
The two-sided bound \eqref{eq:reference-hopping-bounds} yields
\begin{equation}
 \max_{(i,j)\in\mathcal C}|d_{ij}|
   \le C\bigl(m+\log D+\log(1/t_{\min})+1\bigr).
 \label{eq:calibration-increment-bound}
\end{equation}

Apply the constructive linear right inverse in
Lemma~\ref{lem:displacement} to the approximated data.  Its
operator bound is $Cm$ from the maximum edge datum to the maximum
site displacement.  Its formulas use rational operations and
fixed elements of $\mathbb Q(\sqrt3)$, so they can be evaluated
with polynomially many bits.  Round the resulting positions
$Dp_i+s_i$ to rational positions with Euclidean site error at
most $\tau/8$, allocating another such allowance if approximate
arithmetic is used in the right inverse.  Define the final
displacements from these rational positions by
$s_i=X_i-Dp_i$.  This definition preserves the exact contact
vector decomposition even though $s_i$ need not be rational.
Since every $u_{ij}$ is a unit vector, the resulting equations
satisfy
\begin{equation}
 |u_{ij}\cdot(s_j-s_i)-d_{ij}|\le\tau.
 \label{eq:rounded-displacement-equations}
\end{equation}
The right-inverse bound and
\eqref{eq:calibration-increment-bound} imply
\eqref{eq:calibrated-displacement-size}.  For example, the
$\sqrt3$ arithmetic may be done to absolute accuracy
$O(\tau/[D+\poly(m)\max(1,\max|d_{ij}|)])$; the required number
of bits is polynomial.  Noncontacts retain separation at least
$\sqrt2D-2S$ by the triangle inequality.

There is no requirement to realize prescribed distances exactly.
For a contact with $\delta=s_j-s_i$ the actual distance is
\[
 |Du_{ij}+\delta|
   =D+u_{ij}\cdot\delta+O(S^2/D).
\]
Actual coordinates automatically satisfy all vector-cycle
identities; the displayed per-edge correction is already
included in Lemma~\ref{lem:hopping-displacement} and is not
accumulated through a further nonlinear solve.

Put $e_{ij}=u_{ij}\cdot(s_j-s_i)-d_{ij}$.  The key cancellation
is the exact identity
\begin{equation}
 Z_*T_{ij}(X^0)e^{-u_{ij}\cdot(s_j-s_i)}
    =t_{ij}e^{-e_{ij}},\qquad |e_{ij}|\le\tau.
 \label{eq:exact-calibration-identity}
\end{equation}
Lemma~\ref{lem:hopping-displacement} therefore gives
$Z_*T_{ij}(X)=t_{ij}(1+O(\varepsilon_{\mathrm{disp}}+\tau))$.
After tuning, Lemma~\ref{lem:positive-hopping} adds only
$O(\varepsilon_{\mathrm{rel}})$ relative error.  This proves
\eqref{eq:calibrated-contact-error}.  In particular no factor
$e^{Cm}$ or $e^{|d_{ij}|}$ multiplies its absolute error, even
for the smallest floored hopping.  Since $t_{ij}\le1$, the
absolute error per contact is at most $C\mathcal E_{\mathrm{cal}}$;
summing hopping operators incurs only polynomial factors in
the number of sites.

All computations of $X$ used the surrogate fields $W_i^*$.
They depend on the primary centers and on $D$, but not on any
eigenfunction or on the secondary parameters $\rho_i$.  Thus
the algorithm first fixes the rounded charges, computes and
rounds these centers, and only then invokes
Proposition~\ref{prop:tuning} at the final rational centers.
The surrogate-field replacement bounds hold uniformly over the
tuning cube, and the relative orbital comparisons hold uniformly
over all tuned outputs. Thus tuning at these fixed centers requires
no recalibration and introduces no subsequent position-rounding error.

Finally Lemma~\ref{lem:hopping-orientation} and $Z_*\le2e^D$
give $Z_*|B_{ij}-B_{ji}|\le4e^{-0.86D}$.  The common target
energy $-1/2$ is sufficient; exact equality of the $\lambda_i$
is not required.  The full one-body operator is Hermitian, so
one may use either orientation for its off-diagonal entry after
the negligible other-hole terms are accounted for in
Section~\ref{sec:continuum}.
\end{proof}

If the contact set is empty, the procedure simply rounds $Dp_i$
to rational centers at the stated site accuracy; there are no
hopping equations to solve.  In general
\eqref{eq:calibrated-displacement-size} is
$S\le\poly(n)(1+\log D)$.  The bounds needed for the global
choice are precisely \eqref{eq:total-calibration-error}, together
with $m+S\le D^{1/4}$ and the exponentially small orientation
error.  All scale choices and the conversion to the final
many-electron energy error are made in
Section~\ref{sec:completion}.
